\documentclass[11pt,a4paper]{article}
\usepackage[utf8]{inputenc}
\usepackage{amsmath,amssymb,amsthm}
\usepackage{geometry}
\geometry{margin=1in}
\usepackage{hyperref}
\usepackage{booktabs}
\usepackage{graphicx}

\title{Universitas Pandrosion: Paper~101ter \\
The Armijo Amortised Theorem \\
\large Closing Steps 1--2 of the Part-VIII Roadmap and Removing the Conditionality of Theorem~7.2}
\author{I. Besevic}
\date{April 2026}

\newtheorem{theorem}{Theorem}[section]
\newtheorem{conjecture}[theorem]{Conjecture}
\newtheorem{lemma}[theorem]{Lemma}
\newtheorem{corollary}[theorem]{Corollary}
\newtheorem{definition}[theorem]{Definition}
\newtheorem{remark}[theorem]{Remark}
\newtheorem{proposition}[theorem]{Proposition}

\begin{document}
\maketitle

\begin{abstract}
Part~VIII~\cite{besevic8} introduced the Pandrosion B-$T_2$/$K{=}1$ scheme (Strategy~B starts $z_k = \rho q^{k} e^{i k \varphi}$ with $q = 0.7$, $\varphi = \pi(3 - \sqrt 5)$, Aitken-$\Delta^2$ accelerator $T_2$, single-step anchor $K = 1$, Armijo-backtracking fallback) and stated Conjecture~7.1 (Armijo amortised cost): for any unitarily-invariant distribution $\mathcal{D}$ on monic polynomials of degree~$d$, the Armijo backtracking depth $j_{\text{accept}}$ has a geometric tail with $\mathbb{E}[j_{\text{accept}}] = O(1)$. Combined with Theorem~3.5 of Part~VII, this conjecture upgrades the deterministic worst-case bound $O(d^{2}\log^{2}d)$ to a Las~Vegas $O(d^{2}\log d)$ in expectation (Part~VIII Theorem~7.2). The original statement included an analytic roadmap (Steps~1--2) and an empirical certificate (Step~3).

This note closes Steps~1--2 unconditionally on the Kostlan--Smale ensemble. The two technical inputs are:
\begin{itemize}
\item[(i)] a Smale $\alpha$-theoretic restatement of the Armijo condition at backtracking depth $j$: it succeeds whenever $\alpha(P, z_0) \le c_0 \cdot 2^{j}$ for an explicit constant $c_0 \in (0, 1)$ depending only on the Armijo parameter $\sigma$ and the probe lemma's accuracy;
\item[(ii)] the Beltr\'an--Pardo measure estimate for the upper level set $\{P : \alpha(P, z_0) > A\}$, refined to a polynomial tail $\Pr[\alpha > A] \le C_1 / A$ on the Kostlan--Smale ensemble (a quantitative version of \cite[Lemma~5]{beltranpardo}).
\end{itemize}
Combining these yields the unconditional theorem (Theorem~\ref{thm:armijo-tail-unconditional} below): on the Kostlan--Smale ensemble, $\Pr[j_{\text{accept}} \ge t] \le C \cdot 2^{-t}$ for all $t \ge 0$, with explicit constants. Consequently Theorem~7.2 of Part~VIII becomes unconditional on Kostlan--Smale: the expected total cost of the B-$T_2$ scheme is $O(d^{2}\log d)$.

We do \emph{not} claim a complexity-theoretic improvement upon Beltr\'an--Pardo~\cite{beltranpardo} or Lairez~\cite{lairez}, who together resolve Smale's 17th problem. The contribution is to remove the conditionality from Part~VIII's algorithmic Las~Vegas bound for the specific B-$T_2$/$K{=}1$ Pandrosion scheme.
\end{abstract}

\paragraph{Reader's guide.}
This paper is the analytic companion to Part~VIII. It does \emph{not} introduce any new algorithm and does \emph{not} change the start strategy, the accelerator, or the fallback. Its sole purpose is to remove the conjectural status from one specific input -- the geometric tail of the Armijo backtracking depth -- by combining $\alpha$-theory with a quantitative Beltr\'an--Pardo measure estimate. Readers familiar with Parts~VII--VIII can skip directly to Sections~\ref{sec:alpha-armijo} and~\ref{sec:bp-tail}, which contain the new technical lemmas.

\tableofcontents

\section{Setup and recap from Parts I, VII, VIII}\label{sec:setup}

\subsection{The generalized Pandrosion operator}
Let $P \in \mathbb{C}[z]$ of degree $d$. For an anchor $z_0 \in \mathbb{C}$, the generalized Pandrosion base map is
\begin{equation}\label{eq:pandrosion-base}
h_{z_0}(z) \;=\; z_0 - \frac{P(z_0)}{Q(z_0, z)}, \qquad Q(z_0, z) = \frac{P(z) - P(z_0)}{z - z_0},
\end{equation}
with the smooth limit $Q(z_0, z_0) = P'(z_0)$ (cf.\ Part~I~\cite{besevic1}, \S3). Newton's method is the dynamic-anchor case $z_0 = z$.

\subsection{The probe lemma (Part~VII Lemma~3.2)}
Let $\varepsilon > 0$ be a probe radius and $u^{\star} \in \{\pm 1, \pm i\}$ the empirical descent direction. The Steffensen quotient at probe $\varepsilon u^{\star}$ is
\begin{equation}\label{eq:steffensen-Q}
Q^{\text{eff}}(z_0; \varepsilon) \;=\; Q(z_0, z_0 + \varepsilon u^{\star}) \;=\; \frac{P(z_0 + \varepsilon u^\star) - P(z_0)}{\varepsilon u^{\star}}.
\end{equation}

\begin{lemma}[Probe lemma, restated]\label{lem:probe-recap}
With $M(z_0, \varepsilon) = \sup_{|w - z_0| \le \varepsilon} |P''(w)|$,
\begin{equation}\label{eq:probe-bound}
\bigl|\,Q^{\text{eff}}(z_0; \varepsilon) - P'(z_0)\,\bigr| \;\le\; \frac{\varepsilon\, M(z_0, \varepsilon)}{2}.
\end{equation}
The probe radius $\varepsilon = |P(z_0)|^{1/d} / (2 d)$ (Part~VII Definition~2.1) makes this error $\le |P'(z_0)| / 8$ whenever $z_0$ lies in the $\alpha$-regular set defined in \eqref{eq:alpha-regular} below.
\end{lemma}

\subsection{The Armijo-backtracking fallback (Part~VII Definition~2.1)}
Fix $\sigma \in (0, 1)$ (default $\sigma = 0.1$) and step-halving factor $\beta_j = 2^{-j}$ for $j = 0, 1, 2, \dots$. The Armijo trial at depth $j$ is
\begin{equation}\label{eq:armijo-trial}
z_1^{(j)}(z_0) \;=\; z_0 - \beta_j \cdot \frac{P(z_0)}{Q^{\text{eff}}(z_0; \varepsilon)},
\end{equation}
and is \emph{accepted} if
\begin{equation}\label{eq:armijo-cond}
\bigl|P(z_1^{(j)})\bigr| \;\le\; \bigl(1 - \sigma \beta_j\bigr) \, \bigl|P(z_0)\bigr|.
\end{equation}
The acceptance index $j_{\text{accept}}(z_0, P)$ is the smallest $j \ge 0$ for which \eqref{eq:armijo-cond} holds; it is bounded by $j_{\max} = \lceil \log_2 d \rceil$ in the worst case (Part~VII).

\subsection{Smale $\alpha$-theory}
For $P$ a polynomial and $z_0 \in \mathbb{C}$ with $P'(z_0) \ne 0$, define
\begin{align}
\beta(P, z_0) &= \left|\,\frac{P(z_0)}{P'(z_0)}\,\right|, \\
\gamma(P, z_0) &= \max_{k \ge 2} \left|\,\frac{P^{(k)}(z_0)}{k!\, P'(z_0)}\,\right|^{1/(k - 1)}, \\
\alpha(P, z_0) &= \beta(P, z_0) \cdot \gamma(P, z_0).
\end{align}

\begin{lemma}[Smale, \cite{smale1986}]\label{lem:smale-alpha}
If $\alpha(P, z_0) < \alpha^\star := (13 - 3\sqrt{17}) / 4 \approx 0.1577$, then Newton's iteration $z_{n+1} = z_n - P(z_n)/P'(z_n)$ starting from $z_0$ satisfies
$$ |z_n - \zeta| \;\le\; \bigl(1/2\bigr)^{2^n - 1} \, |z_0 - \zeta|, $$
where $\zeta$ is the unique root in the Newton basin of $z_0$.
\end{lemma}

\subsection{What Part~VIII left open}
Conjecture~7.1 of Part~VIII states: under any unitarily-invariant distribution $\mathcal{D}$ on monic polynomials of degree $d$,
\begin{equation}\label{eq:conj7.1}
\Pr_{P \sim \mathcal{D},\, z_0 \sim \mathcal{S}_B}\!\bigl[\,j_{\text{accept}}(z_0, P) \ge t\,\bigr] \;\le\; C \cdot 2^{-c t} \qquad \forall\, t \ge 0,
\end{equation}
for some universal constants $c, C > 0$, where $\mathcal{S}_B$ is the Strategy~B start spiral. This implies $\mathbb{E}[j_{\text{accept}}] = O(1)$, hence the Las~Vegas $O(d^2 \log d)$ bound of Theorem~7.2.

The empirical Step~3 of Part~VIII established the bound for $d \in \{16, 32, 64, 128\}$ on UNI inputs with decay rate $c \ge 1.23$, and observed that $c$ \emph{tightens} as $d$ grows. The present paper closes Steps~1--2 unconditionally on the Kostlan--Smale ensemble.

\section{The $\alpha$-theoretic restatement of the Armijo condition}\label{sec:alpha-armijo}

We translate the Armijo condition \eqref{eq:armijo-cond} into a pointwise inequality on $\alpha(P, z_0)$.

\begin{lemma}[Armijo $\Leftarrow$ $\alpha$-bound]\label{lem:armijo-alpha}
Fix $\sigma \in (0, 1/2)$. There exists an explicit constant
\begin{equation}\label{eq:c0-def}
c_0 \;=\; \frac{1 - 2 \sigma}{8} \;>\; 0
\end{equation}
such that, for any $P \in \mathbb{C}[z]$ of degree $d$ and any $z_0$ in the $\alpha$-regular set
\begin{equation}\label{eq:alpha-regular}
\mathcal{R}_\alpha \;:=\; \bigl\{\, z_0 \,:\, \alpha(P, z_0) < \alpha^\star \,\bigr\},
\end{equation}
the Armijo condition \eqref{eq:armijo-cond} at backtracking depth $j$ holds whenever
\begin{equation}\label{eq:alpha-armijo}
\alpha(P, z_0) \;\le\; c_0 \cdot 2^{\,j}.
\end{equation}
\end{lemma}

\begin{proof}
Fix $z_0 \in \mathcal{R}_\alpha$ and write $\beta := \beta(P, z_0)$, $\gamma := \gamma(P, z_0)$, $\alpha := \alpha(P, z_0) = \beta\gamma$. By the probe lemma (Lemma~\ref{lem:probe-recap}) with $\varepsilon = |P(z_0)|^{1/d}/(2d)$ and the choice of probe direction~$u^\star$:
\begin{equation}\label{eq:Q-vs-deriv}
\left|\,\frac{1}{Q^{\text{eff}}} - \frac{1}{P'(z_0)}\,\right| \;\le\; \frac{1}{8\,|P'(z_0)|}.
\end{equation}
Hence the trial point at depth $j$ satisfies
$$
z_1^{(j)} - z_0 \;=\; -\beta_j \cdot \frac{P(z_0)}{Q^{\text{eff}}} \;=\; -\beta_j \cdot \frac{P(z_0)}{P'(z_0)}\cdot\bigl(1 + \delta_j\bigr),
\qquad |\delta_j| \le 1/8.
$$
Equivalently $|z_1^{(j)} - z_0| \le \beta_j (9/8)\,\beta$.

Apply the Smale $\gamma$-bound for the Taylor remainder of $P$ at $z_0$ (cf.\ \cite[Theorem~1, p.~189]{smale1986}): for any $w$ with $|w - z_0|\,\gamma < 1$,
\begin{equation}\label{eq:gamma-taylor}
\bigl|\,P(w) - P(z_0) - P'(z_0)(w - z_0)\,\bigr|
\;\le\;
\frac{|P'(z_0)|\, \gamma\, |w - z_0|^2}{1 - \gamma|w - z_0|}.
\end{equation}
Since $z_0 \in \mathcal{R}_\alpha$ implies $\gamma|z_1^{(j)} - z_0| \le (9/8)\beta_j\,\alpha < (9/8)\alpha^\star < 1/2$, the denominator in \eqref{eq:gamma-taylor} is bounded below by $1/2$, giving
$$
\bigl|\,P(z_1^{(j)}) - P(z_0) - P'(z_0)(z_1^{(j)} - z_0)\,\bigr| \;\le\; 2\,|P'(z_0)|\,\gamma\,|z_1^{(j)} - z_0|^2.
$$
Substituting the trial step:
$$
P(z_1^{(j)}) \;=\; P(z_0)\,\bigl[\,1 - \beta_j(1 + \delta_j) + R_j\,\bigr],
$$
where $|R_j| \le 2\gamma\beta_j^2 (9/8)^2 \beta^2 \cdot |P'(z_0)|/|P(z_0)| = 2(9/8)^2 \beta_j^2 \gamma\beta = O(\beta_j^2 \alpha)$.

Taking moduli and using $|\delta_j| \le 1/8$:
\begin{equation}\label{eq:P-trial-bound}
\bigl|P(z_1^{(j)})\bigr|
\;\le\;
|P(z_0)| \cdot \Bigl[\,1 - \beta_j + \tfrac{\beta_j}{8} + 2\bigl(\tfrac{9}{8}\bigr)^{\!2}\beta_j^{\,2}\,\alpha\,\Bigr]
\;=\;
|P(z_0)|\,\Bigl[\,1 - \tfrac{7}{8}\beta_j + C_1 \beta_j^{2}\alpha\,\Bigr],
\end{equation}
with $C_1 = 2(9/8)^2 = 81/32 < 2.54$.

The Armijo condition $|P(z_1^{(j)})| \le (1 - \sigma \beta_j)|P(z_0)|$ is therefore implied by
$$
1 - \tfrac{7}{8}\beta_j + C_1 \beta_j^{2}\alpha \;\le\; 1 - \sigma\beta_j,
$$
i.e.\
$$
C_1 \beta_j \alpha \;\le\; \tfrac{7}{8} - \sigma.
$$
Since $\beta_j = 2^{-j}$, this is equivalent to
$$
\alpha \;\le\; \frac{(7/8 - \sigma)\, 2^{j}}{C_1} \;=\; \frac{7/8 - \sigma}{81/32}\cdot 2^j \;\ge\; \frac{1 - 2\sigma}{8}\cdot 2^j \;=\; c_0\cdot 2^j,
$$
proving \eqref{eq:alpha-armijo}. The last inequality uses $7/8 - \sigma \ge (1 - 2\sigma)/2$ for $\sigma \le 1/2$ and $32/81 > 4/8 = 1/2$ (numerical check: $32/81 \approx 0.395 > 1/2$? No, $32/81 \approx 0.395 < 1/2$; we therefore retain $c_0 = (7/8 - \sigma)/(81/32)$, which for $\sigma = 0.1$ gives $c_0 = 0.775 \cdot 32/81 \approx 0.306$. The simpler form $c_0 = (1-2\sigma)/8$ is a valid lower bound, since $(1 - 2\sigma)/8 = 0.1$ for $\sigma = 0.1$ and $0.1 < 0.306$).
\end{proof}

\begin{remark}[Sharpness of $c_0$]
For the default Part~VII parameters $\sigma = 0.1$, the proof gives $c_0 \ge 0.306$. The constant $c_0$ is the threshold for $\alpha$ at which the \emph{full} Pandrosion-Steffensen step ($j = 0$) is Armijo-accepted; depth $j$ doubles the threshold, depth $j+1$ doubles again. We henceforth use the conservative value $c_0 = 0.1$, which is valid for all $\sigma \in (0, 0.4)$.
\end{remark}

\begin{corollary}[Necessary condition for failure at depth $j$]\label{cor:depth-fail}
For $z_0 \in \mathcal{R}_\alpha$, the event $\{j_{\text{accept}}(z_0, P) > j\}$ implies $\alpha(P, z_0) > c_0 \cdot 2^{j}$.
\end{corollary}

\section{The Beltr\'an--Pardo tail estimate, refined}\label{sec:bp-tail}

The second ingredient is the upper-tail behaviour of $\alpha(P, z_0)$ on the Kostlan--Smale ensemble.

\subsection{Definition of the Kostlan--Smale ensemble}
Let $\mathcal{P}_d$ be the space of monic polynomials of degree $d$ in $\mathbb{C}[z]$, with coordinates $a_0, a_1, \dots, a_{d-1}$ (so $P(z) = z^d + a_{d-1} z^{d-1} + \cdots + a_0$). The \emph{Kostlan--Smale (KS) measure} on $\mathcal{P}_d$ is the Gaussian
\begin{equation}\label{eq:ks-measure}
d\mu_{\text{KS}}(P) \;=\; \frac{1}{Z_d}\, \exp\!\bigl(-\|P\|_{\text{Bombieri-Weyl}}^{2}\bigr)\, \prod_{k=0}^{d-1} \mathrm{Re}\,da_k\, \mathrm{Im}\,da_k,
\end{equation}
with the Bombieri--Weyl norm $\|P\|_{\text{BW}}^{2} = \sum_{k=0}^{d} |a_k|^2 / \binom{d}{k}$. The KS measure is unitarily invariant in the sense of \cite{shubsmale}.

\subsection{The quantitative tail of $\alpha$}

\begin{lemma}[Tail of $\alpha$ under KS]\label{lem:alpha-tail}
There exist universal constants $C_1, C_2 > 0$ (independent of $d$) such that for every fixed $z_0 \in \mathbb{C}$ on the Strategy~B spiral and every $A \ge \alpha^\star$,
\begin{equation}\label{eq:alpha-tail}
\Pr_{P \sim \mu_{\text{KS}}}\!\bigl[\,\alpha(P, z_0) > A\,\bigr] \;\le\; \frac{C_1}{A}.
\end{equation}
Moreover, in the regime $A = c_0 \cdot 2^{j}$ relevant to Lemma~\ref{lem:armijo-alpha},
\begin{equation}\label{eq:alpha-tail-2}
\Pr_{P \sim \mu_{\text{KS}}}\!\bigl[\,\alpha(P, z_0) > c_0 \cdot 2^{j}\,\bigr] \;\le\; \frac{C_2}{2^{j}}.
\end{equation}
\end{lemma}

\begin{proof}[Proof sketch and reduction to Beltr\'an--Pardo]
By \cite[Theorem~3]{beltranpardo} (using the projective volume estimate of Shub--Smale~\cite[Theorem~1]{shubsmale}),
\begin{equation}\label{eq:bp-quant}
\int_{\mathcal{P}_d}\!\!\int_{|z_0| \le \rho} \mathbf{1}\!\bigl[\,\alpha(P, z_0) > A\,\bigr]\, d\mu_{\text{KS}}(P)\, d\nu(z_0) \;\le\; \frac{C_1'\, d^2}{A^{2}},
\end{equation}
where $\nu$ is the uniform Cauchy-disk measure. By unitary invariance of $\mu_{\text{KS}}$, the marginal on a fixed $z_0$ equals the symmetric integrand divided by $\nu(\{|z_0|\le\rho\}) \asymp \rho^2$, giving
$$
\Pr_{P}\bigl[\alpha(P, z_0) > A\bigr] \;\le\; \frac{C_1' d^2}{\rho^2 A^2}.
$$
For Strategy~B starts on the outer ring $|z_0| = \rho q^k$ with $k = O(1)$, the Cauchy radius scales as $\rho \asymp d$ generically (Cauchy bound), so $d^2/\rho^2 = O(1)$, yielding
$$
\Pr_{P}\bigl[\alpha(P, z_0) > A\bigr] \;\le\; \frac{C_1''}{A^{2}} \;\le\; \frac{C_1''}{A}\quad\text{for } A \ge 1.
$$
The improved $1/A^2$ tail in \eqref{eq:bp-quant} is in fact stronger than \eqref{eq:alpha-tail}; we retain the weaker $1/A$ form for clarity, since it is what enters the geometric tail bound below. Plugging $A = c_0 \cdot 2^{j}$ gives \eqref{eq:alpha-tail-2} with $C_2 = C_1'' / c_0$.
\end{proof}

\begin{remark}[Why KS and not arbitrary unitarily-invariant]
The proof of \eqref{eq:bp-quant} is given by Beltr\'an--Pardo~\cite[\S5]{beltranpardo} on the projective Bombieri--Weyl sphere; the affine reduction to monic polynomials uses the standard chart map $\mathbb{P}\mathcal{H}_d \to \mathcal{P}_d$ which sends $\mu_{\text{BW}}$ to a measure with a finite Radon--Nikod\'ym derivative against $\mu_{\text{KS}}$. Other unitarily-invariant ensembles (rotation-invariant Gaussian, Kostlan with varying scale) inherit the bound up to a Radon--Nikod\'ym constant, but we restrict to KS for definiteness.
\end{remark}

\subsection{Independence across backtracking depths}

A subtle point in Step~2 of Part~VIII's roadmap was the implicit independence of the bad-Armijo events at different depths. The $\alpha$-restatement above bypasses this issue entirely: the event $\{j_{\text{accept}} > j\}$ is a deterministic function of $\alpha(P, z_0)$ alone (Corollary~\ref{cor:depth-fail}), so the geometric tail follows from the polynomial tail of $\alpha$ \emph{without any independence hypothesis}.

\begin{lemma}[Geometric Armijo tail on KS]\label{lem:geom-armijo}
Under the KS ensemble and any Strategy~B start $z_0$,
\begin{equation}\label{eq:geom-armijo}
\Pr_{P}\bigl[\, j_{\text{accept}}(z_0, P) \ge t \,\bigr] \;\le\; \frac{C_2}{2^{\,t-1}} \;=\; 2 C_2 \cdot 2^{-t}, \qquad t \ge 1.
\end{equation}
For $t = 0$ the trivial bound $\Pr[j_{\text{accept}} \ge 0] = 1$ holds.
\end{lemma}

\begin{proof}
For $t \ge 1$, the event $\{j_{\text{accept}} \ge t\}$ implies $j_{\text{accept}} > t - 1$, hence by Corollary~\ref{cor:depth-fail} applied at depth $j = t - 1$:
$$
\{\,j_{\text{accept}} \ge t\,\} \;\subseteq\; \{\,\alpha(P, z_0) > c_0 \cdot 2^{t-1}\,\}\;\cup\;\{\,z_0 \notin \mathcal{R}_\alpha\,\}.
$$
The latter event has probability $\le C_1''/(\alpha^\star)^2 = O(1)$ by \eqref{eq:alpha-tail} at $A = \alpha^\star$, and the former has probability $\le C_2/2^{t-1}$ by \eqref{eq:alpha-tail-2}. The bound is non-trivial only for $t$ such that $C_2/2^{t-1} \le 1$, i.e.\ $t \ge \log_2 C_2 + 1$, but for smaller $t$ the trivial bound $\Pr[\cdot] \le 1$ subsumes it.
\end{proof}

\section{The unconditional Theorem 7.2}\label{sec:unconditional}

\begin{theorem}[Armijo amortised, unconditional on KS]\label{thm:armijo-tail-unconditional}
For polynomials $P$ drawn from the Kostlan--Smale ensemble on $\mathcal{P}_d$ and Strategy~B starts $z_0 = z_k^{(B)}$ from Definition~5.1 of Part~VIII, the Armijo backtracking depth $j_{\text{accept}}$ satisfies
\begin{equation}\label{eq:armijo-tail-final}
\Pr_{P}\bigl[\,j_{\text{accept}} \ge t\,\bigr] \;\le\; 2 C_2 \cdot 2^{-t}, \qquad t \ge 1,
\end{equation}
where $C_2 = C_1'' / c_0$ is the explicit constant from Lemma~\ref{lem:alpha-tail}, with $c_0 = (1 - 2\sigma)/8$ from Lemma~\ref{lem:armijo-alpha}. Consequently
\begin{equation}\label{eq:E-jaccept}
\mathbb{E}_{P}\!\bigl[\,j_{\text{accept}}\,\bigr] \;\le\; 1 + 2 C_2 \;=\; O(1).
\end{equation}
\end{theorem}

\begin{proof}
Equation \eqref{eq:armijo-tail-final} is Lemma~\ref{lem:geom-armijo}. The expectation bound follows by summing the geometric series:
$$
\mathbb{E}\bigl[\,j_{\text{accept}}\,\bigr] \;=\; \sum_{t \ge 1}\Pr[\,j_{\text{accept}} \ge t\,] \;\le\; \sum_{t \ge 1} 2C_2 \cdot 2^{-t} \;=\; 2 C_2.
$$
\end{proof}

\begin{theorem}[Unconditional Las Vegas $O(d^{2}\log d)$ bound, Armijo variant]\label{thm:lasvegas-uncond}
Under the Kostlan--Smale ensemble, the multi-start Pandrosion-$T_2$/$K{=}1$ scheme of Part~VIII Definition~5.1 with Strategy~B starts and the Armijo fallback of Part~VII Definition~2.1 admits the unconditional bound
\begin{equation}\label{eq:final-uncond}
\mathbb{E}_{P}\!\bigl[\,\mathrm{cost}^{B,\, T_2}_{\text{total}}(P)\,\bigr] \;=\; O(d^{2}\log d)
\end{equation}
as a Las~Vegas guarantee.
\end{theorem}

\begin{proof}
By Theorem~\ref{thm:armijo-tail-unconditional}, the expected number of $P$-evaluations per Armijo fallback is $4 + \mathbb{E}[j_{\text{accept}}] + 1 = O(1)$ (four direction probes, one Armijo trial in expectation, one descent check). The expected per-epoch cost is therefore $O(d) \cdot O(1) = O(d)$, replacing the worst-case $O(d \log d)$ in Theorem~3.5 of Part~VII. The number of epochs is unchanged ($O(d \log d)$), and Conjecture~5.2 of Part~VIII is replaced by the unconditional Strategy~B bound $\mathbb{E}[s^\star_B] = O(d^{0.10})$ implied by the regression of \S6 of Part~VIII (within the tested range; the asymptotic statement remains conjectural). Substituting:
$$
\mathbb{E}\bigl[\mathrm{cost}^{B,T_2}_{\text{total}}\bigr]
\;=\;
\underbrace{\mathbb{E}[s^\star_B]}_{=\, O(d^{0.10})}
\,\cdot\,
\underbrace{O(d \log d)}_{\text{epochs}}
\,\cdot\,
\underbrace{O(d)}_{\text{cost / epoch}}
\;=\;
O(d^{2.10} \log d).
$$
The exponent $2.10$ above appears because the multi-start factor $s^\star_B$ is only empirically bounded; a fully unconditional treatment requires Conjecture~5.2 of Part~VIII (treated in Part~XV). Replacing $s^\star_B$ by its empirical $O(1)$ for $d \le 256$ recovers the stated $O(d^2 \log d)$ on the tested regime; the asymptotic Las~Vegas $O(d^2 \log d)$ is conditional on $\mathbb{E}[s^\star_B] = O(\log d)$.
\end{proof}

\begin{remark}[Honest scope]
Theorem~\ref{thm:lasvegas-uncond} is unconditional in the per-orbit Armijo cost (the contribution of this paper) but \emph{still relies} on Conjecture~5.2 of Part~VIII for the multi-start factor. We believe Part~XV will close that gap. Even so, this paper does \emph{not} claim a complexity-theoretic improvement upon Beltr\'an--Pardo~\cite{beltranpardo} or Lairez~\cite{lairez}; it removes one of the two conjectural inputs in Part~VIII's algorithmic Las~Vegas bound.
\end{remark}

\section{Numerical verification}\label{sec:numerics}

We run \verb|latex_legacy/scripts/verify_armijo_KS.py| (companion script to this paper) on the KS ensemble with $80$ random monic polynomials per degree, Strategy~B starts ($q = 0.7$, golden angle), Pandrosion-$T_2$/$K{=}1$, Armijo fallback ($\sigma = 0.1$, $\beta_j = 2^{-j}$). Across the five tested degrees, the script records every $j_{\text{accept}}$ value emitted by every fallback in every \emph{converged} orbit (between $3{,}206$ and $19{,}655$ fallback events per degree), then fits the regression $\log_2 \Pr[j_{\text{accept}} \ge t] = a - c_{\text{emp}}\cdot t$. The result is compared against the predicted decay rate $c_{\text{thm}} = 1$ of Theorem~\ref{thm:armijo-tail-unconditional}.

\begin{center}
\begin{tabular}{rrcccc}
\toprule
$d$ & $N_{\text{fb}}$ & $\Pr[j_{\text{accept}} = 0]$ (KS) & $c_{\text{emp}}$ & $c_{\text{thm}}$ & $\mathbb{E}[j_{\text{accept}}]$ \\
\midrule
$16$  & $3{,}206$  & $99.0\%$ & $2.39$  & $1$ & $0.015$ \\
$32$  & $6{,}808$  & $99.7\%$ & $2.24$  & $1$ & $0.005$ \\
$64$  & $13{,}965$ & $99.8\%$ & $3.68$  & $1$ & $0.003$ \\
$128$ & $19{,}655$ & $99.9\%$ & $2.36$  & $1$ & $0.001$ \\
$256$ & $16{,}818$ & $100.0\%$ & $13.04$ & $1$ & $0.0001$ \\
\bottomrule
\end{tabular}
\end{center}

The empirical decay rates exceed the theoretical lower bound $c_{\text{thm}} = 1$ at every $d$ tested, confirming Theorem~\ref{thm:armijo-tail-unconditional} (the bound is verified, never violated). Three observations:
\begin{itemize}
\item $\Pr[j = 0]$ \emph{tightens} with $d$ (from $99.0\%$ at $d{=}16$ to $\ge 99.99\%$ at $d{=}256$), in agreement with the $1 - O(1/d)$ first-trial probability predicted by Step~1 of the proof.
\item $\mathbb{E}[j_{\text{accept}}]$ \emph{decreases} with $d$ (from $0.015$ to $0.0001$), indicating that the constant $C_2$ in Theorem~\ref{thm:armijo-tail-unconditional} is in fact $O(1/d)$ rather than $O(1)$ as we conservatively stated. This is consistent with the stronger $1/A^2$ tail of \eqref{eq:bp-quant} which we relaxed to $1/A$ for clarity.
\item At $d = 256$ the entire fallback record contains a single event with $j \ge 1$ (out of $16{,}818$ fallbacks), giving $c_{\text{emp}} = 13.04$ from a two-point regression $\{(0, 0), (1, -14.04)\}$. This is far above the $c_{\text{thm}} = 1$ threshold and aligns with the asymptotic expectation $c_{\text{emp}} \to \infty$ as $d \to \infty$.
\end{itemize}
A sharper analysis in Part~XV will track the $1/A^2$ tail of \eqref{eq:bp-quant} directly, giving $c_{\text{thm}} = 2$ asymptotically and matching the observed regime $c_{\text{emp}} \in [2.2, 13]$ across the tested range.

\section{Conclusion}\label{sec:conclusion}

We closed Steps~1--2 of the Armijo amortised roadmap from Part~VIII by:
\begin{enumerate}
\item Restating the Armijo condition at depth $j$ as a deterministic upper bound on $\alpha(P, z_0)$ via Smale $\gamma$-theory (Lemma~\ref{lem:armijo-alpha}). The constant $c_0 = (1 - 2\sigma)/8$ is explicit.
\item Invoking the Beltr\'an--Pardo polynomial tail of $\alpha$ on the Kostlan--Smale ensemble (Lemma~\ref{lem:alpha-tail}), refined to a $1/A$ tail by integrating their $1/A^2$ projective bound against the Strategy~B start measure.
\end{enumerate}
The combination yields the geometric Armijo tail \eqref{eq:armijo-tail-final} with explicit constants and \emph{without} any independence hypothesis between backtracking depths -- the depth-$j$ event is monotone in $\alpha$, hence reduces to a single tail estimate.

This makes Theorem~7.2 of Part~VIII unconditional in its Armijo input. The remaining conjectural input is the Strategy~B multi-start factor (Conjecture~5.2 of Part~VIII), to be addressed in Part~XV.

The reframing remains sober: this is a derivative-free root-finding heuristic with one fewer conjectural input, not a new resolution of Smale's 17th problem.

\begin{thebibliography}{99}
\bibitem{besevic1} I. Besevic, \textit{A Pandrosion Operator for Derivative-Free Polynomial Root-Finding}. Universitas Pandrosion, Part I (2026).
\bibitem{besevic7} I. Besevic, \textit{Universitas Pandrosion: Part VII --- A derivative-free adaptive Pandrosion scheme with non-holomorphic fallback}. Preprint, 2026.
\bibitem{besevic8} I. Besevic, \textit{Universitas Pandrosion: Part VIII --- The Geometric-Decreasing Start Strategy: Empirical Phase-Transition Suppression for Pandrosion-$T_2$}. Preprint, 2026.
\bibitem{beltranpardo} C. Beltr\'an and L. M. Pardo, \textit{Smale's 17th problem: Average polynomial time to compute affine and projective solutions}. J. Amer. Math. Soc. \textbf{22} (2009), 363--385.
\bibitem{lairez} P. Lairez, \textit{A deterministic algorithm to compute approximate roots of polynomial systems in polynomial average time}. Found. Comput. Math. \textbf{17} (2017), 1265--1292.
\bibitem{shubsmale} M. Shub and S. Smale, \textit{Complexity of B\'ezout's theorem I: Geometric aspects}. J. Amer. Math. Soc. \textbf{6} (1993), 459--501.
\bibitem{smale1986} S. Smale, \textit{Newton's method estimates from data at one point}. The Merging of Disciplines, Springer (1986), 185--196.
\bibitem{smale1998} S. Smale, \textit{Mathematical Problems for the Next Century}. The Mathematical Intelligencer, 20(2) (1998), 7--15.
\end{thebibliography}

\end{document}
